\documentclass[11pt,a4paper]{article}
\usepackage[utf8]{inputenc}
\usepackage[T1]{fontenc}
\usepackage[spanish,es-noquoting]{babel}
\usepackage{amsmath,amssymb}
\usepackage{booktabs}
\usepackage{geometry}
\usepackage{enumitem}
\geometry{margin=2.5cm}

% Generado por correr.py a partir de las soluciones de SCIP. No editar.
\newcommand{\costoOptimo}{2.904.500}
\newcommand{\demandaTotal}{7850}
\newcommand{\stockMaximo}{900}
\newcommand{\costoPuntoTres}{3.128.900}
\newcommand{\propiasPuntoTres}{6530}
\newcommand{\tercPuntoTres}{1320}
\newcommand{\costoPuntoCuatro}{3.128.900}
\newcommand{\propiasPuntoCuatro}{6530}
\newcommand{\tercPuntoCuatro}{1320}
\newcommand{\tablaTercCuatro}{\begin{tabular}{lrrrrrr}
\toprule
Mes & A & B & C & D & E & Total \\
\midrule
Ene/27 & 0 & 0 & 0 & 65 & 0 & 65 \\
May/27 & 20 & 0 & 0 & 0 & 25 & 45 \\
Jun/27 & 0 & 0 & 0 & 0 & 200 & 200 \\
Jul/27 & 75 & 115 & 0 & 0 & 0 & 190 \\
Ago/27 & 200 & 0 & 0 & 0 & 0 & 200 \\
Sep/27 & 40 & 0 & 0 & 0 & 0 & 40 \\
Jun/28 & 75 & 40 & 0 & 45 & 0 & 160 \\
Jul/28 & 20 & 75 & 0 & 0 & 25 & 120 \\
Ago/28 & 0 & 200 & 0 & 0 & 0 & 200 \\
Sep/28 & 0 & 0 & 0 & 0 & 100 & 100 \\
\midrule
Total & 430 & 430 & 0 & 110 & 350 & 1320 \\
\bottomrule
\end{tabular}}
\newcommand{\dminuno}{555}
\newcommand{\costoDminuno}{2.904.500}
\newcommand{\costoNovecientosuno}{2.904.500}
\newcommand{\dmintres}{705}
\newcommand{\costoDmintres}{3.166.300}
\newcommand{\costoNovecientostres}{3.128.900}
\newcommand{\dmincuatro}{705}
\newcommand{\costoDmincuatro}{3.169.300}
\newcommand{\costoNovecientoscuatro}{3.128.900}
\newcommand{\tablaCurvaCinco}{\begin{tabular}{rr}
\toprule
Capacidad $D$ & Costo óptimo \\
\midrule
700 & infactible \\
705 & \$3.169.300 \\
750 & \$3.154.400 \\
800 & \$3.145.900 \\
850 & \$3.137.400 \\
900 & \$3.128.900 \\
1000 & \$3.111.900 \\
1100 & \$3.094.900 \\
1200 & \$3.088.100 \\
\bottomrule
\end{tabular}}


\title{Aplicaciones Computacionales en Negocios\\
\large Trabajo Práctico 2 --- Punto 1: modelo de \emph{lot sizing}}
\author{}
\date{}

\begin{document}
\maketitle

\section{Descripción del problema}

Se debe planificar la producción mensual de $n=5$ productos ($A,\dots,E$) a lo
largo de $T=24$ meses. La demanda
$d_{it}$ de cada producto $i$ en cada mes $t$ se consigue de \texttt{datos.xlsx} y debe cumplirse en su totalidad para cada mes. La
producción de cada producto se realiza en lotes de $L=10$ unidades, con un
máximo de $C=120$ unidades por producto y por mes. El depósito admite, a lo sumo,
$D=900$ unidades de stock entre todos los productos. El stock inicial es nulo y
el costo de producción es $c=\$370$ por unidad, sin costos de almacenamiento.

\section{Modelo de programación lineal entera mixta}

\subsection*{Conjuntos y parámetros}

\begin{align*}
  P &= \{A,B,C,D,E\} && \text{productos,}\\
  T &= \{1,\dots,24\} && \text{meses del horizonte de planificación,}\\
  d_{it} &\ge 0 && \text{demanda del producto } i \text{ en el mes } t,\\
  c &= 370 && \text{costo de producción por unidad (\$/unidad)},\\
  L &= 10 && \text{tamaño del lote (unidades)},\\
  C &= 120 && \text{capacidad de producción por producto y mes (unidades)},\\
  K &= C/L = 12 && \text{cantidad máxima de lotes por producto y mes,}\\
  D &= 900 && \text{capacidad total del depósito (unidades).}
\end{align*}

\subsection*{Variables de decisión}

\begin{align*}
  y_{it} &\in \mathbb{Z}_{\ge 0}, \quad y_{it} \le K
  && \begin{tabular}{l} cantidad de lotes del producto $i$ fabricados en el mes $t$\end{tabular}\\
  s_{it} &\ge 0
  && \begin{tabular}{l} stock del producto $i$ al final del mes $t$ (unidades)\end{tabular}
\end{align*}

La producción en unidades es la expresión $x_{it} = L\,y_{it}$. Modelar la
producción directamente en lotes es lo que garantiza la restricción de
lotificación: al multiplicar una variable entera por $L=10$, $x_{it}$ sólo puede
tomar valores en $\{0,10,20,\dots,120\}$. Las variables de stock se declaran
continuas, por el balance quedan
automáticamente en múltiplos de $5$, ya que las demandas son múltiplos de $5$ y
la producción en múltiplos de $10$.

\subsection*{Formulación}

\begin{align}
  \min\quad & z = \sum_{i\in P}\sum_{t\in T} c\,L\,y_{it}
  \label{eq:obj}\\[2pt]
  \text{s.a.}\quad
  & s_{i,t-1} + L\,y_{it} - d_{it} = s_{it},
  && \forall i \in P,\ \forall t \in T
  \label{eq:balance}\\
  & s_{i0} = 0,
  && \forall i \in P
  \label{eq:inicial}\\
  & \sum_{i \in P} s_{it} \le D,
  && \forall t \in T
  \label{eq:deposito}\\
  & 0 \le y_{it} \le K, \quad y_{it} \in \mathbb{Z},
  && \forall i \in P,\ \forall t \in T
  \label{eq:capacidad}\\
  & s_{it} \ge 0,
  && \forall i \in P,\ \forall t \in T
  \label{eq:nonneg}
\end{align}

\begin{description}[leftmargin=1.6cm,style=nextline,itemsep=2pt]
  \item[\eqref{eq:obj}] Minimiza el costo total de producción: \$370 por cada
    unidad fabricada, es decir \$3.700 por lote.
  \item[\eqref{eq:balance}] Balance de inventario: el stock con el que se cierra
    un mes es el stock del mes anterior más lo producido menos lo demandado.
    Junto con $s_{it}\ge 0$ garantiza que toda la demanda se cumpla en el mes en
    que se produce (no se admiten faltantes ni demanda atrasada).
  \item[\eqref{eq:inicial}] Stock inicial nulo. No tenemos nada a priori de ningún producto.
  \item[\eqref{eq:deposito}] Capacidad del depósito: el stock total almacenado
    al final de cada mes, sumado sobre todos los productos, no supera las 900
    unidades.
  \item[\eqref{eq:capacidad}] Capacidad de producción y lotificación: a lo sumo
    12 lotes (120 unidades) de cada producto por mes, en cantidades enteras de
    lotes.
  \item [\eqref{eq:nonneg}] No negatividad de cualquier producto en cualquier mes. Trivial.
\end{description}

\section{Resolución}

El modelo se implementó en Python con la biblioteca \texttt{PuLP} y se resolvió
con el solver CBC (archivo \texttt{punto1.py}, la demanda se lee de
\texttt{datos.xlsx}). El solver alcanza el óptimo:

\[
  z^{\star} = \$\,\costoOptimo
\]

\section{Punto 2: límite de 300 unidades por mes}

El modelo es el mismo del punto 1 más una única familia de restricciones, que
acota la producción total mensual sumada sobre todos los productos:
\begin{equation}
  \sum_{i \in P} L\,y_{it} \le 300, \qquad \forall t \in T
  \label{eq:limite-mensual}
\end{equation}
Son 24 restricciones nuevas; no hacen falta variables adicionales. Notar que
\eqref{eq:limite-mensual} es más fuerte que la capacidad individual
\eqref{eq:capacidad}: antes se podían fabricar hasta $5\times120=600$ unidades
por mes.

\paragraph{¿Se puede cumplir con la demanda total?} No: el modelo es
\textbf{infactible} (así lo reporta el solver, archivo \texttt{punto2.py}). La
razón se ve sin necesidad de resolverlo. Como el stock inicial es nulo, la
producción acumulada hasta el mes $t$ no puede superar $300\,t$, y debe alcanzar
para cubrir la demanda acumulada:
\[
  \sum_{\tau \le t} \sum_{i \in P} d_{i\tau} \;\le\; 300\,t,
  \qquad \forall t \in T .
\]
La condición se viola por primera vez en agosto de 2027, donde la demanda
acumulada es de 2820 unidades contra una capacidad acumulada de 2400 (faltan 420
unidades), y vuelve a violarse desde julio de 2028 en adelante. Sobre el
horizonte completo, la demanda total de \demandaTotal{} unidades excede la
capacidad total de $300 \times 24 = 7200$ unidades: \textbf{faltan 650
unidades}. Con este límite es imposible satisfacer toda la demanda, de modo que
el modelo sólo tiene sentido si se relaja el problema --- permitiendo demanda
insatisfecha (con un costo de faltante) o agregando capacidad externa, que es
justamente lo que introduce el punto 3 con la tercerización.

\end{document}
